% 中文节标题：结论
\FloatBarrier
\section{Conclusion}

% 中文：本文研究跨镜头切换的流式重建问题，其中递归历史具有相互冲突的作用：它为建立相邻镜头之间的关系提供必要参照，
% 但持续传播同一状态又可能在新镜头内引起随时间变化的漂移。我们的核心观察是，历史应被临时读取以完成镜头间对齐，
% 而不应作为新镜头的递归状态继续传播。基于这一观察，Shot3R通过切换前历史推断镜头间关系，同时仅允许独立重新初始化的状态
% 驱动后续重建。它将各镜头中的相机、人体和场景几何表示到统一世界坐标系，并维持跨镜头人物身份，同时按照输入顺序处理每一帧，
% 无需访问未来帧或进行逐视频优化。三个多人体基准上的实验表明，Shot3R能够持续改善统一坐标系下的人体重建、相机轨迹估计和
% 人物身份关联，并在大幅视角变化下取得最明显的收益。如何在包含复杂剪辑、人物反复进出以及快速视觉变化的长视频中保持稳定重建，
% 仍是值得进一步研究的问题。
We study streaming reconstruction across shot transitions, where recurrent
history plays conflicting roles: it provides the reference needed to relate
successive shots, yet its continued propagation can introduce time-varying
drift within the new shot. Our key insight is that history should be
temporarily read for inter-shot alignment without being propagated as the
recurrent state of the new shot. Based on this insight, \method{} infers the
inter-shot relation from pre-transition history while allowing only an
independently reinitialized state to drive subsequent reconstruction. It
expresses the cameras, humans, and scene geometry from each shot in the common
world frame and maintains person identities while processing frames in
arrival order, without future-frame access or per-video optimization.
Experiments on three multi-person benchmarks demonstrate consistent
improvements in world-frame human reconstruction, camera trajectory
estimation, and identity association, with the largest gains under substantial
viewpoint changes. Handling long videos with complex edits, repeated entrances
and exits, and rapid visual changes remains an important direction for future
work.
